\subsection{SUMMABLE FAMILIES IN A COMMUTATIVE GROUP}

In this section we shall be concerned only with Hausdorff commutative topological groups, whose law of composition is written additively. Only the most important results will be translated into multiplicative notation.

Let \(\mathbf{G}\) be a Hausdorff commutative group, let \(\mathbf{I}\) be any index set and let \((x_{i})_{i\in \mathbf{I}}\) be a family of points of \(\mathbf{G}\) , indexed by \(\mathbf{I}\) . With each finite subset \(\mathbf{J}\) of \(\mathbf{I}\) we associate the element \(s_{\mathbf{J}} = \sum_{i\in \mathbf{J}}x_i\) of \(\mathbf{G}\) , which

we call the finite partial sum of the family \((x_{i})_{i\in I}\) , corresponding to the set \(J\) . If \(\mathfrak{F}(I)\) denotes the set of finite subsets of \(I\) , we have thus a mapping \(J\to s_J\) of \(\mathfrak{F}(I)\) into \(G\) . Now \(\mathfrak{F}(I)\) is a directed set with respect to the relation \(\subset\) ; for if \(J\) and \(J'\) are two elements of \(\mathfrak{F}(I)\) , then \(J\subset J\cup J'\) and \(J'\subset J\cup J'\) , and \(J\cup J'\) is a finite subset of \(I\) . Let \(\Phi\) be the section filter of the directed set \(\mathfrak{F}(I)\) (Chapter I, § 6, no. 3).

DEFINITION 1. Let \((x_{i})_{i\in I}\) be a family of points of a Hausdorff commutative group G; let \(\mathfrak{F}(I)\) be the set of finite subsets of the index set \(I\) , and for each finite subset \(J\) of \(I\) , let \(s_J\) be the sum of those \(x_{i}\) such that \(i\in J\) . The family \((x_{i})_{i\in I}\) is said to be summable if the mapping \(J\to s_J\) has a limit with respect to the section filter \(\Phi\) of the set \(\mathfrak{F}(I)\) directed by the relation \(\subset\) ; this limit is then said to be the sum of the family \((x_{i})_{i\in I}\) and is denoted by \(\sum_{i\in I}x_{i}\) (or simply \(\sum x_{i}\) , or even \(\sum x_{i}\) , when there is no risk of ambiguity).

Definition 1 is equivalent to the following: the family \((x_{i})\) is summable and its sum is \(s\) if, for each neighbourhood \(\mathbf{V}\) of the origin in \(\mathbf{G}\) , there is a finite subset \(\mathbf{J}_0\) of \(\mathbf{I}\) such that for each finite subset \(\mathbf{J} \supset \mathbf{J}_0\) of \(\mathbf{I}\) we have \(s_{\mathbf{J}} \in s + \mathbf{V}\) .

If G is written multiplicatively, and if we put \(p_{J} = \prod_{i \in J} x_i\) for each finite subset J of I, the family \((x_i)\) will be said to be multipliable if the mapping \(J \to p_J\) has a limit with respect to the filter \(\Phi\) ; this limit is called the product of the family \((x_i)\) , and is denoted by \(\prod_{i \in I} x_i\) .

Remarks. 1) When I is finite, Definition 1 reduces to the ordinary definition of the sum of a finite family. More generally, if I is arbitrary and \(x_{i} = 0\) except when the index \(i\) belongs to a finite subset J of I, then the sum \(\sum_{i \in I} x_{i}\) is equal to \(\sum_{i \in J} x_{i}\) .

\begin{enumerate}
\def\labelenumi{\arabic{enumi})}
\setcounter{enumi}{1}
\item
  The definition of a summable family does not involve any ordering of the index set I, and we may therefore say that the notion of a sum thus defined is commutative. More precisely, we have the following property: let \((x_{t})_{t\in \mathbf{I}}\) be a summable family and let \(\varphi\) be a bijective mapping of an index set K onto the set I; then if we put \(y_{x} = x_{\varphi (x)}\) , the family \((y_{x})_{x\in \mathbb{K}}\) is summable and has the same sum as \((x_{t})\) . For if \(s = \sum_{i\in I}x_i\) , and if \(\sum_{i\in J}x_i\in s + V\) for every finite subset J containing the finite subset \(J_0\) , then we shall have \(\sum_{x\in L}y_x\in s + V\) for every finite subset L of K containing \(\varphi^{-1}(J_0)\) .
\item
  Definition 1 applies, more generally, to every family of points in a Hausdorff topological space X on which has been defined an associative and commutative law of composition, written additively; for it does not make use of the axioms of topological groups.
\end{enumerate}

